\section{The Standard Model on forty points}
\label{sec:physics}

Sections 2--8 established exact structure. This section asks the question the programme
exists to answer: how much of the actual physics of our universe does that structure
carry? We give the answer at three levels --- what the finite geometry itself contains,
what a string compactification built on it produces, and what is still missing --- and we
end with a scorecard.

\begin{figure}[htbp]
\centering
\begin{tikzpicture}[
   pill/.style={rounded corners=5pt,font=\sffamily\scriptsize\bfseries,text=white,minimum width=23mm,minimum height=5mm,inner sep=2pt},
   row/.style={anchor=west,font=\small}]
  \def\entry#1#2#3#4{%
    \node[pill,fill=#3] at (1.2,-#1*0.62) {#4};
    \node[row] at (2.75,-#1*0.62) {#2};}
  \entry{0}{Quantum measurement structure (qutrit Paulis, contexts, Clifford group)}{teal}{EXACT}
  \entry{1}{Exceptional algebras $E_6$, $E_7$, $E_8$, $F_4$ and the Albert algebra}{teal}{EXACT}
  \entry{2}{A finite light cone; clocks; one Lorentzian real form $\so(1,9)$}{teal}{EXACT}
  \entry{3}{Gauge group $SU(3)\times SU(2)\times U(1)$}{sky}{HOSTED}
  \entry{4}{Three generations}{sky}{HOSTED}
  \entry{5}{One chiral generation as a spinor ($16$ of $\mathrm{Spin}(1,9)$)}{teal}{EXACT}
  \entry{6}{Which chirality: exchanged only by antiunitary (CPT-type) symmetries}{rose}{NO-GO}
  \entry{7}{Matter parity in a supersymmetric singlet vacuum (heterotic class)}{rose}{EXCLUDED}
  \entry{8}{A supersymmetric, realistic vacuum (nonabelian directions)}{gold}{NOT YET}
  \entry{9}{Masses, mixing angles, coupling constants}{slate}{OPEN}
  \entry{10}{Gravity, a dynamical space-time, the cosmological constant}{slate}{OPEN}
\end{tikzpicture}
\caption{Scorecard. \textsf{EXACT}: a theorem or exact computation in this paper.
\textsf{HOSTED}: realised in an explicit construction built on the geometry, not forced by
it. \textsf{EXCLUDED}: exhaustively ruled out in a stated scope (Section~\ref{sec:sneutrino}).
\textsf{NO-GO}: proved impossible from inside. \textsf{OPEN}: not addressed by any exact result.}
\label{fig:scorecard}
\end{figure}

\subsection{What the finite geometry itself contains}

\paragraph{The anatomy of one symmetry.}
Pick a point $p_0$ of $\W$ and the \emph{transvection} $R$ centred on it: the symplectic
map $x\mapsto x+\langle x,v\rangle v$, of order three. It fixes the $13$ points of the plane
$p_0^\perp$ --- $p_0$ and its $12$ neighbours --- and permutes the other $27$ points in nine
orbits of three.

\begin{figure}[htbp]
\centering
\begin{tikzpicture}[scale=1,>={Stealth[length=1.6mm]}]
  \fill[ink] (0,0) circle (3pt); \node[font=\scriptsize,anchor=north] at (0,-0.1) {$p_0$};
  \foreach \k in {0,...,11}{\fill[teal] ({30*\k+15}:1.05) circle (2pt);}
  \draw[teal!30] (0,0) circle (1.05);
  \foreach \t in {0,...,8}{
     \pgfmathsetmacro{\a}{40*\t}
     \foreach \j in {0,1,2}{
        \coordinate (q\t\j) at ({\a+(\j-1)*11}:2.35);}
     \draw[->,rose!70,line width=0.6pt] (q\t0) to[bend left=35] (q\t1);
     \draw[->,rose!70,line width=0.6pt] (q\t1) to[bend left=35] (q\t2);
     \draw[->,rose!70,line width=0.6pt] (q\t2) to[bend right=50] (q\t0);
     \foreach \j in {0,1,2}{\fill[rose] (q\t\j) circle (2pt);}}
  \node[font=\sffamily\scriptsize,align=left,anchor=west] at (3.1,1.0)
   {\textcolor{teal!80!black}{\textbf{fixed:}} $p_0$ and its $12$ neighbours};
  \node[font=\sffamily\scriptsize,align=left,anchor=west] at (3.1,0.4)
   {\ \ centraliser $C(R)$, order $648$, acts on the\\\ \ $12$ with rank $3$: \ $\C^{12}=\mathbf1\oplus\mathbf3\oplus\mathbf8$};
  \node[font=\sffamily\scriptsize,align=left,anchor=west] at (3.1,-0.55)
   {\textcolor{rose!80!black}{\textbf{moved:}} the $27$ far points, in nine $3$-cycles};
  \node[font=\sffamily\scriptsize,align=left,anchor=west] at (3.1,-1.15)
   {\ \ eigenspaces of $R$ on $\C^{27}$: \ $9\oplus9\oplus9$};
  \node[font=\sffamily\scriptsize,align=left,anchor=west] at (3.1,-1.75)
   {\ \ centre of $C(R)$ $=\langle R\rangle\cong\Z_3$};
\end{tikzpicture}
\caption{One transvection of $\W$. Its centraliser acts on the twelve fixed neighbours
with the dimensions of the Standard Model's gauge algebra, $1+3+8$, and it splits the
twenty-seven moved points into three equal eigenspaces.}
\label{fig:transvection}
\end{figure}

\begin{result}{Proposition 9.1 (the transvection spine) \hfill \tierC}
The centraliser $C(R)$ of a transvection has order $648$ (it is the point stabiliser
$3^{1+2}{:}2A_4$) and centre $\langle R\rangle\cong\Z_3$. Its permutation module on the $12$
neighbours of $p_0$ decomposes as $\mathbf1\oplus\mathbf3\oplus\mathbf8$. On the $27$ far
points, $R$ has eigenspaces $9\oplus9\oplus9$.
\cert{bt886\_standard\_model\_spine.py}
\end{result}

\begin{boundary}[How far this goes]
The dimensions $1,3,8$ are those of the Lie algebras $\mathfrak u(1)$, $\mathfrak{su}(2)$,
$\mathfrak{su}(3)$, and the corpus has sometimes read Proposition~9.1 as ``the Standard
Model gauge group follows with zero free parameters''. That is an over-read. The
$\mathbf1\oplus\mathbf3\oplus\mathbf8$ here are representations of a \emph{finite} group of
order $648$; they are not Lie algebras, they do not come with gauge fields, and nothing
fixes their couplings. The honest statement is that the local symmetry of one point of the
geometry has the \emph{shape} of the Standard Model's gauge algebra. \tierB
\end{boundary}

\paragraph{Three from the Steinberg module.}
The first homology of $\W$ (its $40$ points, $240$ edges and $160$ triangles in lines) is
$81$-dimensional and irreducible under $\PSp(4,3)$: it is the \emph{Steinberg module}, the
most rigid representation a group of Lie type has, and the logical space of the
repository's $[[240,81,3]]_3$ quantum code. \tierC\ (\cert{bt861\_code\_register\_is\_steinberg.py})

\begin{result}{Theorem 9.2 (every order-three symmetry splits matter $27+27+27$) \hfill \tierT}
Let $g\in\PSp(4,3)$ have order three. Then $g$ acts on the $81$-dimensional Steinberg
module with eigenvalues $1,\omega,\omega^2$, each with multiplicity exactly $27$.
\end{result}

\begin{proof}
The Steinberg character vanishes on every element whose order is divisible by the
defining characteristic $p=3$ (a theorem of Steinberg; see \cite{Humphreys}). So $\chi(g)=\chi(g^2)=0$, and the
multiplicity of each eigenvalue $\omega^j$ is
$\tfrac13\bigl(\chi(1)+\omega^{-j}\chi(g)+\omega^{-2j}\chi(g^2)\bigr)=81/3=27$.
\end{proof}

\begin{boundary}[Three is forced, but not selected]
Theorem~9.2 holds for \emph{every} element of order three, so it produces a threefold
splitting without choosing one. And the string-theory track (below) found that in the
explicit compactifications the number of families is the number of complex planes of the
internal torus, not a property of $\W$. We therefore record the Steinberg splitting as a
structural host for three generations, not a derivation of them. \tierB
\end{boundary}

\paragraph{The qutrit register is $SU(9)$, not $E_6$.}
One might hope that the $E_6$ of grand unification acts directly on qutrit registers.

\begin{result}{Theorem 9.3 (no Pauli $E_6$) \hfill \tierT\ \tierC}
Conjugation by the three-qutrit Pauli group splits $\mathfrak{gl}(27)$ into $729$ distinct
one-dimensional character spaces, so every conjugation-stable subalgebra is spanned by
Pauli operators. A simple Pauli-spanned Lie algebra has dimension $3^a-3^b$, i.e.\ one of
$8,24,72,80,216,240,648,720,728$. Since $\dim\e_6=78$ is not among them, $E_6$ cannot act
on a qutrit register through Pauli-stable generators.
\cert{Holotrade 1cde1f1}
\end{result}

The natural gauge algebra of the qutrit register is instead $\sll_9$ --- the $80$ of
Section~5 --- which is exactly what appears in the string construction.

\subsection{The string that carries the geometry}
\label{sec:strings}

The most developed attempt to turn this structure into physics uses the
$E_8\times E_8$ heterotic string compactified on an orbifold \cite{DHVW}: six extra dimensions curled
up as three two-dimensional tori, divided by a finite rotation group. Its advantage is
that everything is computable exactly, with a public tool, the \textsc{orbifolder}
\cite{Orbifolder}.

\begin{result}{Theorem 9.4 ($\W$ lives in the twisted sector of $E_8$) \hfill \tierC}
For the $\Z_3$ twist of the $E_8$ lattice vertex algebra, the commutator form of the
twisted-sector Pauli group equals twice the lattice form on every rung, and the resulting
commutation geometry is $\W$. The twist is anomalous on a single $E_8$ (ground-state energy
$\rho=4/9$), consistent exactly when $3\mid k$ on $E_8^{\,k}$; the $E_8^3$ orbifold is the
Niemeier lattice vertex algebra $V_{A_8^3}$ with algebra $\sll_9^{\,3}$.
\cert{Holotrade 68df33f}
\end{result}

\begin{figure}[htbp]
\centering
\begin{tikzpicture}[>={Stealth[length=2.3mm]},
   st/.style={draw=#1,fill=#1!10,rounded corners=4pt,align=center,font=\sffamily\scriptsize,minimum width=31mm,minimum height=13mm},
   ar/.style={->,thick,ink!55}]
  \node[st=ink] (a) at (0,0) {$E_8\times E_8$ heterotic string\\\scriptsize the twist of $\W$\\\scriptsize $=$ an $A_8$ shift};
  \node[st=rose] (b) at (4.3,0) {order $3$: $T^6/\Z_3$\\\scriptsize $SU(9)\times SO(14)\times U(1)$\\\scriptsize \textbf{no complete family}};
  \node[st=teal] (c) at (8.6,0) {order $6$: $T^6/\Z_6$\\\scriptsize $42$ of $79$ SMs from\\\scriptsize the $\W$ vacuum};
  \node[st=gold] (d) at (8.6,-2.3) {flagship model\\\scriptsize $SU(3)\times SU(2)\times U(1)_Y$,\\\scriptsize three net families};
  \node[st=slate] (e) at (4.3,-2.3) {the vacuum problem\\\scriptsize SUSY, $\mu$, proton,\\\scriptsize neutrinos};
  \draw[ar] (a) -- (b); \draw[ar] (b) -- node[above,font=\tiny]{even order} (c);
  \draw[ar] (c) -- (d); \draw[ar] (d) -- (e);
  % three tori
  \begin{scope}[xshift=0cm,yshift=-2.3cm]
    \foreach \i/\t/\c in {0/{\tfrac16}/sky,1/{\tfrac16}/sky,2/{-\tfrac13}/gold}{
       \begin{scope}[xshift=\i*1.25cm-1.25cm]
       \draw[\c,line width=1pt,fill=\c!10] (0,0) ellipse (0.55 and 0.33);
       \draw[\c,line width=0.8pt] (-0.25,0.03) .. controls (-0.1,-0.09) and (0.1,-0.09) .. (0.25,0.03);
       \draw[\c,line width=0.8pt] (-0.17,-0.03) .. controls (-0.05,0.05) and (0.05,0.05) .. (0.17,-0.03);
       \node[font=\scriptsize] at (0,-0.55) {$\t$};
       \end{scope}}
    \node[font=\sffamily\scriptsize,text=gold!70!black,align=center] at (1.25,0.6) {qutrit plane};
    \node[font=\sffamily\scriptsize,text=slate] at (0,-1.0) {twists of the three tori};
  \end{scope}
\end{tikzpicture}
\caption{The heterotic route. At order three no Wilson-line breaking completes a
Standard-Model family (a theorem); at even order the $\W$ vacuum is the most productive
source of Standard Models; the open problem is a realistic vacuum. The third torus, the
unique order-three plane, is where the Higgs lives.}
\label{fig:heterotic}
\end{figure}

\paragraph{Order three fails, for a reason.}
In the $T^6/\Z_3$ orbifold built on this twist, Wilson lines produce three-family $SU(5)$
and $SO(10)$ models; but an order-three Wilson line keeps at most one of $Q,u^c,e^c$ in each
$\mathbf{10}$, so no breaking completes a Standard-Model family. The general law: for a
$\Z_N$ orbifold, $u^c$ and $e^c$ survive together only when $N$ is even. This explains, from
first principles, why heterotic model builders moved to even-order orbifolds. An exhaustive
census of $127\,980$ classes confirms it. \tierT\ \tierC\ (Holotrade \cert{ed1eaca},
\cert{a6e1c69}, \cert{a779053})

\paragraph{Order six succeeds.}
Among the $58$ $\Z_6$-I and $61$ $\Z_6$-II shifts, the $\W$ vacuum
$SU(9)\times SO(14)\times U(1)$ is the most productive class: of $79$ inequivalent Standard
Models found in it, $42$ come from the $\Z_6$-I $\W$ vacuum. The flagship has gauge group
$SU(3)_C\times SU(2)_L\times U(1)_Y$ with three net families and no anomalies, and at one
fixed point the local gauge group is $SU(9)$ with all four simple roots of the
Standard Model's $SU(5)$ inside it: the Standard Model's unification group descends from the
twist's own $\sll_9$. \tierC\ (Holotrade \cert{5b3f3ad})

\begin{result}{Measured regularities of the $\W$ Standard Models \hfill \tierC}
\begin{itemize}[leftmargin=1.2em]
\item \textbf{The breaking is an entangling gate.} The order-three Wilson line's $\sll_9$
      holonomy is of two-qutrit controlled-$Z$ type $(5,2,2)$ in $45$ of $45$ $\Z_6$-I models,
      against a $33\%$ background (chance $\approx2.6\times10^{-22}$).
      (\cert{d772153})
\item \textbf{Two commuting holonomies cut out $SU(3)\times SU(2)$}: in all $25$ models with
      parity-type $\theta^3$, against a $30\%$ background ($\approx8.6\times10^{-14}$).
      (\cert{94df376})
\item \textbf{A heavy top is allowed}: $57$ of $60$ models admit a renormalisable top Yukawa
      coupling. (\cert{93b34e1})
\item \textbf{The Higgs sits alone in the qutrit plane}: in all $87$ $\Z_6$-I models the weak
      doublets come only from the third torus (the unique order-three plane) and colour
      triplets only from the other two; this one fact gives both the top coupling and
      \textbf{doublet--triplet splitting}, solved in $55$ of $87$ models by a
      missing-partner mechanism. (\cert{3caf15e}, \cert{a6f1cae})
\end{itemize}
\end{result}

\begin{physical}[What it means physically]
The Standard Model's most puzzling structural feature --- that the Higgs doublet is light
while its would-be colour-triplet partner, which would make protons decay, is absent --- is
here a consequence of geometry: the Higgs lives on the one internal plane whose twist has
order three, the ``qutrit plane'', and that plane carries no colour triplet. And the
symmetry breaking that produces the Standard Model is, in qutrit language, an entangling
gate. These are measured regularities of an explicit string construction, with controls;
they are not theorems about all constructions.
\end{physical}

\subsection{Where the string route stands}

A model is not a universe until it has a stable vacuum with the right low-energy
properties. Here the honest record is mixed, and it has been corrected more than once.

\begin{center}
\small
\renewcommand{\arraystretch}{1.3}
\begin{tabularx}{\textwidth}{@{}>{\raggedright\arraybackslash}p{3.0cm}X>{\raggedright\arraybackslash}p{2.4cm}@{}}
\toprule
\textbf{Requirement} & \textbf{Current exact status} & \textbf{Badge}\\
\midrule
Supersymmetric vacuum & In the doublet--triplet-solving vacuum, $F$- and $D$-flatness conflict:
one $U(1)$ direction gives every mass singlet a positive charge. Gordan's alternative
(``$D$-flat or empty superpotential, never both'') verified on all $215$ models. & \tierX\ (that vacuum)\\
$D$-flat directions & Exist in $23$ of $87$ ($\Z_6$-I) and $105$ of $128$ ($\Z_6$-II); $83$ of
the $211$ anomalous models have none. & \tierC\\
Proton stability & No supersymmetric singlet vacuum preserves a matter parity built from the
gauge torus and the space group: $0$ of $215$ (Section~\ref{sec:sneutrino}). Baryon triality is
impossible in all $215$. Heavy superpartners would need $m_t\le171$\,GeV. & \tierC\ (excluded)\\
Parity from rotations & Opens in one model ($\Z_6$II$_{23}$): unbroken, the exotics stay massless;
broken, they are heavy but no Higgs pair stays light. & \tierC\ / \tierO\\
The $\mu$ term & Arises at order $3$ or $4$ in every model; none is $\mu$-free through order
five. Must come from a vacuum symmetry. & \tierO\\
Neutrino masses & In the splitting-solving models, one to two orders of magnitude too light. & \tierO\\
\bottomrule
\end{tabularx}
\end{center}

Figure~\ref{fig:gauntlet} gathers the whole search in one picture. Each band is an exact filter; the
number is how many models or vacua pass it, and the tag names the computation that decided it.

\begin{figure}[htbp]
\centering
\begin{tikzpicture}[
  band/.style={draw=#1!70!black,fill=#1soft,rounded corners=2pt,minimum height=0.78cm,
               font=\sffamily\scriptsize,align=center,line width=0.6pt},
  tag/.style={font=\sffamily\tiny,text=slate,anchor=west,align=left},
  num/.style={font=\sffamily\bfseries\small,text=ink}]
  % Z6 column
  \node[font=\sffamily\bfseries\small,text=ink] at (0,0.9) {$\Z_6$ orbifolds};
  \node[band=sky,minimum width=7.2cm] (a) at (0,0) {Standard Models with the $W(3,3)$ twist \quad \textbf{215}};
  \node[band=sky,minimum width=5.6cm] (b) at (0,-1.05) {a $B-L$ matter parity exists \quad \textbf{88}};
  \node[band=sky,minimum width=4.2cm] (c) at (0,-2.1) {$D$-flat with FI cancelled \quad \textbf{24}};
  \node[band=rose,minimum width=5.6cm] (d) at (0,-3.15) {parity survives (torus $\times$ space group) \quad \textbf{0}};
  \node[band=gold,minimum width=5.6cm] (e) at (0,-4.2) {parity from plane rotations \quad \textbf{1} model};
  \node[band=gold,text width=3.0cm] (f) at (-1.95,-5.45) {38 minimal vacua:\\ exotics massless};
  \node[band=gold,text width=3.0cm] (g) at (1.95,-5.45) {larger vacua: exotics heavy;\\ flat only by cancellation};
  \foreach \u/\v in {a/b,b/c,c/d,d/e} \draw[->,ink!50,line width=0.7pt] (\u) -- (\v);
  \draw[->,ink!50,line width=0.7pt] (e) -- (f); \draw[->,ink!50,line width=0.7pt] (e) -- (g);
  \node[tag] at (3.75,0) {5b3f3ad};
  \node[tag] at (2.95,-1.05) {P10960};
  \node[tag] at (2.25,-2.1) {P10960};
  \node[tag] at (2.95,-3.15) {P10960--10967};
  \node[tag] at (2.95,-4.2) {P10968};
  % side results
  \node[band=rose,minimum width=3.6cm] (t) at (6.4,-1.05) {baryon triality \quad \textbf{0}/215};
  \node[tag] at (4.65,-1.55) {P10965 ($SU(5)$ relation)};
  \node[band=rose,minimum width=3.6cm] (h) at (6.4,-2.6) {heavy superpartners:\\ need $m_t\le171$\,GeV};
  \node[tag] at (4.65,-3.3) {P10969 (Higgs quartic)};
  % Z12 column
  \node[font=\sffamily\bfseries\small,text=ink] at (6.4,-4.05) {$\Z_{12}$-I orbifolds};
  \node[band=sky,minimum width=3.6cm] (z1) at (6.4,-4.75) {Standard Models \quad \textbf{289}};
  \node[band=sky,minimum width=3.0cm] (z2) at (6.4,-5.6) {viable parity \quad \textbf{32}};
  \node[band=sky,minimum width=2.6cm] (z3) at (6.4,-6.45) {with a vacuum \quad \textbf{14}};
  \node[band=rose,minimum width=2.6cm] (z4) at (6.4,-7.3) {exotics massive \quad \textbf{0}};
  \foreach \u/\v in {z1/z2,z2/z3,z3/z4} \draw[->,ink!50,line width=0.7pt] (\u) -- (\v);
  % Z2 x Z6 column
  \node[font=\sffamily\bfseries\small,text=ink] at (0,-6.55) {$\Z_2\times\Z_6$-I orbifolds};
  \node[band=sky,minimum width=3.2cm] (y1) at (-2.55,-7.3) {Standard Models \quad \textbf{29}};
  \node[band=sky,minimum width=2.3cm] (y2) at (0.45,-7.3) {any parity \quad \textbf{2}};
  \node[band=rose,minimum width=2.3cm] (y3) at (3.2,-7.3) {survives FI \quad \textbf{0}};
  \foreach \u/\v in {y1/y2,y2/y3} \draw[->,ink!50,line width=0.7pt] (\u) -- (\v);
\end{tikzpicture}
\caption[The vacuum gauntlet]{The vacuum gauntlet. Blue: filters passed; red: dead ends; gold: the one
door still ajar. Every count is exact (Farkas certificates, Smith normal forms, lattice tests and
integer programs), and each is regenerated by the script named in the ledger.}
\label{fig:gauntlet}
\end{figure}

\subsection{The Fayet--Iliopoulos term forces a right-handed sneutrino}
\label{sec:sneutrino}

Every model of the class has an anomalous $U(1)$, whose Fayet--Iliopoulos term must be
cancelled by giving vacuum values to some singlet fields ($D$-flatness), with $\langle
n\rangle/M_s\approx0.2$--$0.4$. Proton stability needs a \emph{matter parity} --- a $\Z_2$
under which quarks and leptons are odd and Higgs fields even --- to survive that
condensation, so the fields that condense must all be even. The earlier record found every
tested $D$-flat direction contained an odd singlet, but it had \emph{sampled} $60$ choices of
$B-L$ per model and said explicitly that this was not a proof; the raw spectra needed to
finish it were in no repository. They were recovered for this paper from the original scan
and frozen as an exact-rational ledger of all $215$ models.

\begin{result}{Theorem 9.6 (no matter parity survives the FI term) \hfill \tierC}
Over the \emph{entire} rational solution set $B-L=x_0+Nt$ of each model:
\begin{enumerate}[label=(\roman*),leftmargin=1.5em]
\item $64$ of the $88$ models with a $B-L$ direction have no $D$-flat singlet direction at all
      (an exact Farkas vector certifies that the FI term cannot be cancelled);
\item in the other $23$ $\Z_6$-I models, the singlets that are even for \emph{some} choice of
      $B-L$ are never $D$-flat (one exact Farkas vector each); the only singlets that are
      never even carry $3(B-L)=3$, i.e.\ $B-L=+1$, fixed by the Standard Model itself, and every
      $D$-flat direction uses one;
\item in the one remaining ($\Z_6$-II) model, all $304$ FI-cancelling extreme rays of the $D$-flat
      cone ($3334$ rays, computed exactly) are obstructed by an integer (Smith normal form)
      parity certificate; its never-even singlets carry $B-L=\pm1$.
\end{enumerate}
Independently, all $1024$ $\Z_2$ characters of the charge lattice that are odd on every
$q,u^c,d^c,e^c$ give no even $D$-flat direction. Controls: the same exact machinery finds
the $23+105=128$ $D$-flat models and the $88$ $B-L$ models of the earlier record.
\cert{w33\_pass10960\_matter\_even\_dflat\_closure.py}
\end{result}

\begin{figure}[htbp]
\centering
\begin{tikzpicture}[x=0.052cm,y=1cm,
   seg/.style={draw=white,line width=0.8pt},
   lab/.style={font=\sffamily\scriptsize,text=white}]
  % row 1: 215 models
  \node[anchor=east,font=\sffamily\scriptsize] at (-3,0) {$215$ models};
  \fill[slate!55] (0,-0.3) rectangle (4,0.3);
  \fill[rose!65] (4,-0.3) rectangle (87,0.3);
  \fill[teal!60] (87,-0.3) rectangle (215,0.3);
  \node[lab] at (45.5,0) {83: FI term cannot be cancelled};
  \node[lab] at (151,0) {128: $D$-flat directions exist};
  \node[font=\tiny,text=slate,anchor=south] at (2,0.32) {4 anomaly-free};
  % row 2: 88 with B-L
  \node[anchor=east,font=\sffamily\scriptsize] at (-3,-1.1) {$88$ with $B-L$};
  \fill[rose!65] (0,-1.4) rectangle (64,-0.8);
  \fill[gold!75] (64,-1.4) rectangle (87,-0.8);
  \fill[ink!80] (87,-1.4) rectangle (88,-0.8);
  \node[lab] at (32,-1.1) {64: no $D$-flat direction};
  \node[font=\sffamily\scriptsize,anchor=west] at (90,-1.1) {23 + 1: every $D$-flat vacuum condenses a $B-L=\pm1$ field};
  % row 3: Z2 characters
  \node[anchor=east,font=\sffamily\scriptsize] at (-3,-2.2) {$1024$ $\Z_2$ parities};
  \fill[rose!65] (0,-2.5) rectangle (160,-1.9);
  \node[lab] at (80,-2.2) {1024: no even $D$-flat direction};
  % verdict
  \node[draw=rose,fill=rosesoft,rounded corners=3pt,font=\sffamily\small\bfseries,text=rose!85!black]
       at (107.5,-3.3) {matter-even supersymmetric singlet vacua: \ $0$};
\end{tikzpicture}
\caption{The exhaustive closure of Theorem~9.6. Every one of the $215$ models is accounted for;
no branch leaves a matter-parity-preserving supersymmetric vacuum in singlet directions.}
\label{fig:closure}
\end{figure}

\begin{physical}[What it means physically: the sneutrino that breaks R-parity]
A singlet with $B-L=+1$ has the quantum numbers of a right-handed (s)neutrino. Theorem~9.6
says the anomalous $U(1)$ \emph{forces} such a field to condense. Martin's criterion
\cite{Martin1992} is exactly relevant here: matter parity survives the breaking of $B-L$ only
if the fields that break it have even $3(B-L)$ --- which is why realistic $SO(10)$ models
break $B-L$ with $\mathbf{126}$-type fields ($B-L=\pm2$) rather than $\mathbf{16}$-type ones
($B-L=\pm1$). In this class the FI term leaves no choice: $B-L$ is broken by one unit, and
the baryon- and lepton-number-violating couplings that matter parity was protecting are no
longer forbidden by any symmetry, at the scale of the singlet vacuum values.
\end{physical}

\paragraph{Why: one $E_6$-shaped charge.}
The obstruction has a single mechanism. In each of the $23$ models there is an exact
certificate $f$ --- a combination of the anomalous and ordinary $U(1)$ charges --- with a
prescribed physical shape: $f=\tfrac{24}{5}$ on every $q,u^c,e^c$ and $f=\tfrac85$ on every
$d^c,\ell$ (in units of a positive scale), and $f\ge0$ on every singlet that could be even.
On the singlets $f$ takes only the values $-8$, $0$ and $8$, and the negative ones are
exactly the forced $B-L=+1$ fields. Because $D$-flatness makes the $f$-charge-weighted sum of
condensates equal the (negative) FI contribution, one of them must condense. On the families
$f=4Q_\psi-\tfrac45Q_\chi$, the $U(1)$ combination of
$E_6\supset SO(10)\times U(1)_\psi\supset SU(5)\times U(1)_\chi\times U(1)_\psi$ with
$Q_\chi=4Y-5(B-L)$; the value $-8$ at $(B-L,Y)=(1,0)$ then identifies the forced field as the
right-handed-sneutrino direction of the matter-odd $\mathbf{16}_{-3}$ in the adjoint $\mathbf{78}$ of
$E_6$. \tierC\ (\cert{w33\_pass10964\_e6\_shaped\_fi\_mechanism.py}) \ \tierB\ (the $E_6$ reading)

\begin{figure}[htbp]
\centering
\begin{tikzpicture}[x=0.62cm,>={Stealth[length=2mm]}]
  \draw[->,ink!60,line width=0.8pt] (-11,0) -- (10.5,0) node[right,font=\scriptsize]{$f$};
  \foreach \v in {-10,-8,-6,-4,-2,0,2,4,6,8}{\draw[ink!40] (\v,-0.08) -- (\v,0.08);
     \node[font=\tiny,text=slate,anchor=north] at (\v,-0.1) {$\v$};}
  \fill[rosesoft] (-11,0.15) rectangle (0,1.9);
  \node[font=\sffamily\tiny,text=rose!80!black,anchor=north west] at (-10.9,1.88) {FI side: needs $f<0$};
  % singlets
  \foreach \v/\c/\l in {-8/rose/{forced $\nu^c$\\$\mathbf{16}_{-3}\subset\mathbf{78}$},0/slate/{singlets},8/sage/{singlets}}{
     \fill[\c] (\v,0.55) circle (4pt);
     \node[font=\sffamily\tiny,align=center,anchor=south] at (\v,0.75) {\l};}
  % families
  \fill[teal] (4.8,-0.95) circle (3pt); \node[font=\sffamily\tiny,anchor=north] at (4.8,-1.1) {$q,u^c,e^c$};
  \fill[teal] (1.6,-0.95) circle (3pt); \node[font=\sffamily\tiny,anchor=north] at (1.6,-1.1) {$d^c,\ell$};
  % higgs
  \fill[gold] (-9.6,-0.95) circle (3pt); \node[font=\sffamily\tiny,anchor=north] at (-9.6,-1.1) {$H_u\in\mathbf{10}_{-2}$};
  \fill[gold] (-1.6,-0.95) circle (3pt); \node[font=\sffamily\tiny,anchor=north] at (-1.6,-1.1) {$H_u$ (other)};
  \node[font=\sffamily\scriptsize,text=slate] at (0,-2.1) {singlets above the axis, charged fields below (units of $\mu$; all $23$ models)};
\end{tikzpicture}
\caption{The canonical certificate $f=4Q_\psi-\tfrac45Q_\chi$ on the families. Among singlets,
only the right-handed-sneutrino direction of the $E_6$ adjoint's $\mathbf{16}_{-3}$ sits on the
side that can balance the Fayet--Iliopoulos term.}
\label{fig:fline}
\end{figure}

\paragraph{The hidden sector does not rescue it.}
Theorem~9.6 covers singlets of every nonabelian group. Fields that are Standard-Model
singlets but charged under the \emph{hidden} gauge group can also condense, and their parity
can be compensated by a hidden gauge element. Pass~10962 attacks this loophole in tiers:
\begin{itemize}[leftmargin=1.2em]
\item \textbf{exact, for every hidden direction:} a parity-preserving vacuum must be $D$-flat for
      the hidden maximal torus with component-wise parity (up to a hidden torus element);
      for $32$ of the $88$ models even this relaxation has an exact Farkas certificate;
\item \textbf{exact over explicit invariant generators:} by the Luty--Taylor and Kempf--Ness
      correspondence a parity-preserving $D$-flat vacuum exists iff an FI-cancelling product of
      matter-even hidden-group invariants does; over singlets, all quadratic invariants,
      $SU(N)$ baryons and $SU(5)$ cubics, the other $55$ models close by exact Farkas vectors,
      and the last by a parity obstruction;
\item the four anomaly-free models (no FI term) admit no matter parity at all.
\end{itemize}
\tierC\ (\cert{w33\_pass10962\_hidden\_sector\_dflat\_tiers.py})

\begin{boundary}[Scope]
For the $56$ models not closed at the first tier, completeness of the invariant generators is
not proved (they contain bi-fundamental, $SU(4)$~$\mathbf6$ or $SU(5)$~$\mathbf{10}$ hidden
fields), so for them the hidden-sector statement is ``no rescue among the standard
invariants'' rather than a theorem. Whether higher-order selection rules suppress the
regenerated couplings in a particular model is not addressed. \tierO
\end{boundary}

\paragraph{The other discrete guardian is excluded too.}
Matter parity is not the only discrete symmetry that can keep the proton alive. \emph{Baryon
triality} $B_3=\exp\bigl(2\pi i(B-2Y)/3\bigr)$ (Ib\'a\~nez--Ross) forbids the baryon-number
violating couplings $u^cd^cd^c$ and $qqql$ while allowing the lepton-number violating ones.
Could the class carry it instead? As a gauge symmetry it would have to be a $\mathbb{Z}_3$
character of the $U(1)$ charge lattice, equal to $B_3$ on every quark and lepton up to a
hypercharge rotation. Pass~10965 shows that this linear system over $\mathrm{GF}(3)$ is
inconsistent in all $215$ models, and names the reason: in $214$ of them single copies obey
the exact identity
\[
  Q(u^c)+Q(e^c)=2\,Q(q)
\]
of $U(1)$ charge vectors --- the extra $U(1)$s act on the $SU(5)$ $\mathbf{10}$ through the
$\mathbf{10}$ itself. On that combination $B_3$ is $1+1-0=2$ while every hypercharge shift is
$2+0-2=0$ mod $3$, so nothing repairs it; the remaining model has an explicit integer
relation with the same property. The zero-target control is solvable in all $215$. This is
the structural reason behind Holotrade's observation that the three dimension-four operators
are locked together (b81ef8c). \tierC\ (\cert{w33\_pass10965\_baryon\_triality\_impossible.py})

\begin{plainwords}
Of the two textbook discrete symmetries that protect the proton, the first dies at the
Fayet--Iliopoulos term and the second is ruled out by the grand-unified shape of the
charges. What remains is an $R$-symmetry --- the anomaly-free $\mathbb{Z}_4^R$ of Lee, Raby,
Ratz, Ross, Schieren, Schmidt-Hoberg and Vaudrevange --- which in a string orbifold lives in the
internal rotation sector rather than in the gauge lattice. That is the one open route.
\end{plainwords}

\paragraph{The vacuum switches the forbidden couplings back on.}
Two gaps remained. First, the searches above used the $B-L$ family and $\mathbb{Z}_2$ characters
that are $\pm1$ on \emph{every} field. A matter parity only has to be $-1$ on quarks and leptons
and $+1$ on the condensate; exotic fields may carry any phase. Second, the string has discrete
symmetries of its own, the \emph{space-group selection rules}, and these are not part of the gauge
lattice. Pass~10967 closes both. It takes the most general candidate: any element of the gauge
torus times the space-group group (the point-group $\mathbb{Z}_6$, a $\mathbb{Z}_3$, and in
$\mathbb{Z}_6$-II a further $\mathbb{Z}_2\times\mathbb{Z}_2$ from fixed points). It then decides
exactly, with a Smith normal form, which singlets such an element can leave even. The answer is
unchanged: such parities exist in the same $88$ models and survive the FI term in none.

The reason is a single identity. In $290$ of the $302$ forced singlets, single copies of the fields
satisfy
\[
  Q(n)\;=\;-Q(u^cd^cd^c)\;=\;-Q(q\,\ell\,d^c)\;=\;-Q(\ell\,\ell\,e^c),
\]
which is the $SO(10)$ quartic invariant $\mathbf{16}^4$ evaluated at $\nu^c$. Each cubic is odd
under every matter parity, so $n$ is too. The string's own coupling rules then finish the argument
(Figure~\ref{fig:rpv}). In each of the $23$ $D$-flat $\mathbb{Z}_6$-I models, none of the three
$R$-parity-violating cubics is allowed at the renormalisable level. Yet all nine forced singlets
of every model appear in allowed quartics $n\,u^cd^cd^c$, $n\,q\,\ell\,d^c$ and $n\,\ell\,\ell\,e^c$.
The vacuum that supersymmetry requires turns all three back on, with strength
$\langle n\rangle/M_s\approx0.2$--$0.4$. Holotrade had noted that its class closure (b81ef8c)
held only if the vacuum breaks matter parity; this is that vacuum, named field by field.
\tierC\ (\cert{w33\_pass10967\_fi\_vacuum\_regenerates\_rpv.py})

\begin{figure}[htbp]
\centering
\begin{tikzpicture}[>={Stealth[length=2mm]},
  bx/.style={draw=#1,fill=#1soft,rounded corners=3pt,font=\sffamily\scriptsize,align=center,
             minimum height=1.25cm,text width=3.0cm,line width=0.7pt}]
  \node[bx=gold,text width=2.6cm] (fi) at (0,0) {FI term $\xi\neq0$\\ must be cancelled};
  \node[bx=rose] (n) at (4.4,0) {forced singlet $n$\\ $Q(n)=-Q(u^cd^cd^c)$\\ odd under every parity};
  \node[bx=sky] (q) at (9.4,0) {string-allowed quartics\\ $n\,u^cd^cd^c$, $n\,q\ell d^c$, $n\,\ell\ell e^c$};
  \node[bx=rose,text width=4.2cm] (r) at (4.4,-2.3) {all three RPV couplings, coefficient $\sim\langle n\rangle/M_s$\\ proton decays far too fast};
  \draw[->,ink!70,line width=0.8pt] (fi) -- node[above,font=\sffamily\tiny]{$\langle n\rangle\neq0$} (n);
  \draw[->,ink!70,line width=0.8pt] (n) -- node[above,font=\sffamily\tiny]{23/23 models} (q);
  \draw[->,ink!70,line width=0.8pt] (q.south) |- (r.east);
  \draw[->,ink!70,line width=0.8pt] (n.south) -- (r.north);
\end{tikzpicture}
\caption[Why supersymmetric vacua violate $R$-parity]{Why supersymmetric vacua of the class violate $R$-parity. The only fields that can cancel
the Fayet--Iliopoulos term carry exactly the conjugate charge of the three dimension-four
operators. Every string selection rule allows the resulting quartic couplings, so the vacuum turns
the cubic operators back on.}
\label{fig:rpv}
\end{figure}

\paragraph{Could heavy superpartners make it harmless?}
Another part of this programme argues that the geometry's supersymmetry is \emph{structural}, with no
light superpartners. Heavy squarks do tame $R$-parity violation, because proton decay through
$\lambda'\lambda''$ falls as $1/\tilde m^4$. With $\lambda'\sim\lambda''\sim\langle n\rangle/M_s\approx0.2$--$0.4$,
the bound $|\lambda'\lambda''|\lesssim10^{-27}(\tilde m/100\,\mathrm{GeV})^2$ needs
$\tilde m\gtrsim6\times10^{14}$--$1.2\times10^{15}$\,GeV. The Higgs sets a ceiling. Running the
Standard-Model couplings up at two loops (checked against the published three-loop value at the
Planck scale), the Higgs self-coupling turns negative at $3.2\times10^{10}$\,GeV (Figure~\ref{fig:lambda}).
Supersymmetric matching requires it to be non-negative wherever the superpartners appear. Pushing
that scale to $10^{15}$\,GeV would need a top quark of at most $171.0$\,GeV, $5.4$ standard deviations
below the measured $172.57\pm0.29$\,GeV. \tierC\ (\cert{w33\_pass10969\_rpv\_heavy\_superpartners\_vs\_higgs\_quartic.py})

\begin{figure}[htbp]
\centering
\begin{tikzpicture}[x=0.52cm,y=55cm,>={Stealth[length=2mm]}]
  \fill[rose!15] (14.78,-0.018) rectangle (19.09,0.135);
  \draw[->,ink!60] (2,0) -- (19.6,0) node[right,font=\scriptsize]{$\log_{10}\mu/\mathrm{GeV}$};
  \draw[->,ink!60] (2,-0.018) -- (2,0.135) node[above,font=\scriptsize]{$\lambda(\mu)$};
  \foreach \x in {4,8,12,16}{\draw[ink!40] (\x,0.001)--(\x,-0.001); \node[font=\tiny,anchor=north,text=slate] at (\x,-0.001) {$\x$};}
  \foreach \y in {-0.01,0.04,0.08,0.12}{\draw[ink!40] (2.1,\y)--(1.9,\y); \node[font=\tiny,anchor=east,text=slate] at (1.9,\y) {$\y$};}
  \node[font=\sffamily\tiny,text=rose!80!black,align=center] at (16.95,0.11) {squarks needed\\ to tame RPV};
  \draw[teal,line width=1.1pt] plot[smooth] coordinates {(2.237,0.1262) (2.500,0.1149) (3.000,0.0967) (3.500,0.0817) (4.000,0.0692) (4.500,0.0585) (5.000,0.0493) (5.500,0.0414) (6.000,0.0345) (6.500,0.0285) (7.000,0.0232) (7.500,0.0185) (8.000,0.0144) (8.500,0.0107) (9.000,0.0075) (9.500,0.0047) (10.000,0.0022) (10.500,0.0000) (11.000,-0.0019) (11.500,-0.0035) (12.000,-0.0050) (12.500,-0.0063) (13.000,-0.0073) (13.500,-0.0083) (14.000,-0.0090) (14.500,-0.0097) (15.000,-0.0102) (15.500,-0.0106) (16.000,-0.0109) (16.500,-0.0111) (17.000,-0.0112) (17.500,-0.0112) (18.000,-0.0112) (18.500,-0.0111) (19.000,-0.0110) (19.086,-0.0109)};
  \fill[rose] (10.51,0) circle (2pt) node[above left,font=\sffamily\tiny,text=rose!80!black]{$3.2\times10^{10}$\,GeV};
  \node[font=\sffamily\tiny,text=teal!80!black,anchor=west] at (5.2,0.075) {measured $m_t=172.57$\,GeV};
\end{tikzpicture}
\caption{The Higgs self-coupling run to high energy (two loops, NNLO matching). It turns negative
near $3\times10^{10}$\,GeV, far below the $\gtrsim6\times10^{14}$\,GeV that heavy squarks would need to
make the regenerated $R$-parity violation harmless (shaded).}
\label{fig:lambda}
\end{figure}

\paragraph{The last door: a parity made of rotations.}
The discrete groups tested so far live in the gauge lattice and the space-group rules. The orbifold
also has \emph{rotational} symmetries, turning each internal plane by its own angle, and these act on
the supersymmetry charges ($R$-symmetries). Combinations of them that leave the superpotential alone
can act as a matter parity. This reopens exactly one model, $\mathbb{Z}_6$II$_{23}$. There, $38$ minimal
supersymmetric vacua keep such a parity, with the right Higgs content; the parity is an honest symmetry
of all $1524$ couplings the string allows. But the unbroken rotation that makes these vacua flat is
the same one that forbids the exotic colour triplets from getting masses, at every order.

Larger vacua break that rotation. For the first complete parity assignment the search finds a vacuum
of $68$ directions (singlets and composites of a hidden $SU(5)$), all kept even by \emph{one} parity
element, $D$-flat with the Fayet--Iliopoulos term cancelled, in which every exotic mass matrix
reaches full rank by fifth order (exotics near $10^{15}$--$10^{16}$\,GeV), and no singlet outside the
vacuum can spoil flatness, because parity forbids it. What it cannot do is keep a Higgs pair light:
the Higgs mass matrix is full rank too. So the trade has two sides, and the class lands on one of
them in every case checked:
\begin{center}
\small
\begin{tabular}{lccc}
\toprule
vacuum & flat by symmetry & exotics & a light Higgs pair \\
\midrule
rotation unbroken ($202$ found) & yes & massless & protected \\
rotation broken, $68$ directions & no & heavy & not protected \\
rotation partly broken, $23$ directions & no & heavy & protected \\
\bottomrule
\end{tabular}
\end{center}
The last row is the \emph{approximate} rotation symmetry that the problem called for: broken at low
order for the exotics and never for one Higgs pair. It is certified exactly (a positive $D$-flat vector, one
parity element for the whole vacuum, every mass term checked against every selection rule, and every
direction outside the vacuum barred from spoiling flatness). What it lacks is flatness \emph{by
symmetry}: its superpotential is not zero, starting at sixth order with a term built from two mesons
of the hidden $SU(5)$. The branch search that makes the superpotential vanish always lands back on
the first row, with the exotic $d$-quarks massless --- $202$ times out of $202$. Whether the hidden
$SU(5)$, whose dynamics is famously able to generate a superpotential of its own, supplies the
missing cancellation is the sharpest open question the programme now has. The same analysis on a
third family of orbifolds, $\Z_2\times\Z_6$ with the $W(3,3)$ twist, finds $29$ Standard Models and
no surviving parity. \tierC\ (\cert{w33\_pass10968\_r\_symmetry\_parity\_and\_massless\_exotics.py})

\subsection{Handedness: host, not select}
\label{sec:chirality}

The weak force acts only on left-handed particles. Can $\W$ explain why?

\begin{result}{Theorem 9.5 (the chirality no-go) \hfill \tierT\ \tierC}
The two half-spin representations $S^\pm$ of the $D_5$ structure built on the
geometry are exchanged by the substrate's own outer controller $T$, which generates
$\PGSp(4,3)=\PSp(4,3){:}2=W(E_6)$ together with $\PSp(4,3)$ and has $\det T=-1$. Consequently
no invariant of the full automorphism group of $\W$ can distinguish $S^+$ from $S^-$.
\cert{w33\_pass346\_the\_chirality\_no\_go.py}
\end{result}

\begin{plainwords}
The forty-point geometry is perfectly capable of \emph{carrying} a left-handed world: the
chiral spinors are there (Section~8). What it cannot do is \emph{choose} left over right,
because one of its own symmetries swaps them. Any explanation of the handedness of the weak
force must come from outside the geometry --- from a choice of vacuum, of orientation, or of
dynamics. This is the same distinction as in Section~2: the two groups of order $51\,840$
differ precisely by whether that swapping symmetry is present.
\end{plainwords}

\begin{physical}[What it means physically: the swap is antiunitary]
Which symmetries do the swapping? For the qutrit, the determinant of a symplectic label map
decides: determinant $+1$ is a unitary Clifford operation, determinant $-1$ can only be
implemented antiunitarily (Pass~5730,
\cert{w33\_pass5727\_5730\_torsion\_family\_heisenberg.py}). The decisive statement was
proved on the $32$ orientations of a complete two-qutrit factorisation frame (Holotrade
\cert{e92a047}, \cert{frame\_chirality\_is\_reversed\_only\_by\_antiunitaries.py}): they split
$16+16$ exactly as the two half-spin weight sets of $D_5$, unitary Cliffords preserve each
half, and antiunitary ones --- and nothing else --- exchange them. The other track has since
found the same law on the clock: determinant-odd ticks exchange the two $D_4$ half-spin
sheets (Pass~10955) and are not completely positive (Pass~10952). So the geometry's
chirality is invariant under every \emph{unitary} symmetry and exchanged only by
\emph{antiunitary} ones. That is the structure of a chiral quantum theory obeying CPT:
the mirror of a chiral world is related to it by an antiunitary map, and which of the two is
called ``left'' is a convention that no internal measurement can fix. The no-go is the
finite form of this, not a defect peculiar to $\W$. \tierB\
(\cert{w33\_pass10952\_clock\_complete\_positivity\_firewall.py},
\cert{w33\_pass10955\_d4\_halfspin\_clock\_bridge.py})
\end{physical}

\begin{result}{Corollary 9.7 (one coset reverses time, acts antiunitarily, and swaps chirality) \hfill \tierT}
Fix a ``now'' (the Bell line $B$). Its stabiliser in $\PGSp(4,3)$ has order $1296$, and
its elements outside $\PSp(4,3)$ form a single coset of size $648$. That coset is exactly
\begin{enumerate}[label=(\roman*),leftmargin=1.5em]
\item the set of symmetries that \emph{reverse} the unique global orientation of the
      history graph (Theorem~4.4);
\item the set that can only be implemented \emph{antiunitarily} on qutrits (Pass~5730);
\item the set that \emph{exchange} the two chiral half-spin modules (Theorem~9.5;
      Holotrade \cert{e92a047}).
\end{enumerate}
\end{result}

\begin{proof}
Each item identifies its set with the complement of $\PSp(4,3)$ in the relevant group:
(i) is the statement of Theorem~4.4(ii) (the $648$ elements of the Bell stabiliser in
$\PSp(4,3)$ fix the orientation cycle, the other $648$ negate it); (ii) symplectic similitudes of
multiplier $+1$ lift to unitary Clifford operations and those of multiplier $-1$ only to
antiunitary ones; (iii) $\PSp(4,3)$ preserves $S^+$ and $S^-$ (they are inequivalent) while the
outer element $T$ exchanges them. Intersecting with the stabiliser of $B$ gives the same coset
in all three cases.
\end{proof}

\begin{physical}[What it means physically: a finite CPT theorem]
In quantum field theory, the CPT theorem says that the only way to relate a chiral world to
its mirror image by a symmetry is antiunitary, and that the same operation reverses the
direction of time. Corollary~9.7 is that structure, exactly, in the forty-point geometry: the
operations that reverse time's orientation are the antiunitary ones, and they are precisely
the ones that exchange left and right. Choosing an orientation of time does not choose a
chirality --- but it makes chirality a conserved label, and the world of opposite chirality
is the time-reversed one. This answers, in the finite model, the question of whether the
datum that fixes the arrow and the datum that fixes handedness are related: they live in the
same coset. \tierB
\end{physical}

\subsection{What we do not count as evidence}

The older manuscripts contain many numerical coincidences: $1/\alpha\approx137$ written as a
combination of cyclotomic values of $q=3$, spectral indices of the cosmic microwave
background, spanning-tree counts read as gravitational entropies. Some are striking. None
is derived from a dynamical principle with its uncertainties, and a finite geometry with
many natural integers can match many targets. We therefore do not use any of them as
evidence in this paper. (The value $\sin^2\theta_W=3/8$ at unification, often quoted
alongside them, is the standard prediction of any $SU(5)$ or $SO(10)$ grand unification and
is not specific to $\W$.) \tierB\ at most.

\begin{keyidea}
The geometry carries the Standard Model's shapes --- its gauge algebra's dimensions, three
equal matter sectors, a chiral spinor, a matter parity --- and a string compactification
built on its twist produces genuine three-family Standard Models whose Higgs sits in the
``qutrit plane''. It cannot choose the handedness of the weak force --- only an antiunitary,
CPT-type symmetry exchanges the two --- and, exhaustively, none of its supersymmetric singlet
vacua preserves a matter parity: the anomalous $U(1)$ forces a $B-L=\pm1$ condensate.
\end{keyidea}
